% Formen-Rechner in Plain TeX. TODO: Seitenzahlen
\magnification=\magstep1
\hsize=15cm \vsize=22cm
\font\gross=cmbx12 scaled \magstep2
\newcount\anzahl \anzahl=3
\newdimen\rand \rand=1.5in
\def\titel#1{\bigskip\centerline{\gross #1}\medskip}
\titel{Der Formen-Rechner}
\noindent Ein {\bf Kreis} hat die Fl\"ache $A = \pi r^2$ und den Umfang
$$U = 2\pi r \eqno(1)$$
\beginsection 1. Formen

\item{$\bullet$} Kreis \hfill {\it rund}
\item{$\bullet$} Rechteck \quad {\tt a mal b}
\loop \ifnum\anzahl>0 \advance\anzahl by -1 \repeat
\ifx\undefined\x \message{fehlt} \else \relax \fi
\halign{#\hfil&\quad#\cr Form&Wert\cr Kreis&3.14\cr}
\hbox to \hsize{\hss rechts}\vskip 12pt plus 2pt minus 1pt
\input formen
\bye
